\documentclass[paper=b5j]{jlreq}

\usepackage{hyperref}
\usepackage[margin=20truemm]{geometry}

\title{還元主義はままならない}
\author{Stanaka}

\begin{document}
\maketitle
\setlength{\parskip}{0.3\baselineskip}


某放課後、篠澤広とそのプロデューサーの事務所。

「…ふふん、さらり。…ふふ、さらっ」

チラチラとこちらを見ながら、わざとらしく髪をかき上げる広さんに訝しげな視線を向ける。

「広さん、そういうのは口に出すものではありませんよ。…切っては…いませんね」

「ふふ、どう？」

手の甲で髪を持ち上げたまま、得意げな表情の広さん。


「だから、切ったわけではないでしょう？」

「うん。でも、どう？」    

「どうと言われましても...。いつも通りかわいいですよ」

「ふふ。…じゃなくて、いつも通り？」

それでは納得がいかないとばかりに、身を乗り出してくる。

「…こほん。いつも以上にかわいいです」

「そうなんだ。やっぱり星南はすごい」

満足げに眼を細めながら、その目線はこの場には居ない友人に向いているようだ。

「十王星南さんがどうかしましたか？」

「星南がね、私の『あいどるぱわー』がまた上がった、って」

彼女はそう言って立ち上がり、胸に手を当てポーズを取って見せる。思わず顎に手を当て、広さんを凝視してはみるが…結局は分からず、首をかしげてしまった。

「前々から思っていましたが、いったい彼女には何が見えているのでしょうね」

「おじいちゃんにも見えるって。わたし、そのおかげで入学できた」

「ふむ。学園長や星南さんにどのようにやるのか聞いてみるのもいいかもしれませんね。プロデュースの役に立つかもしれませんし」

人差し指を立て提案してみるが、手を下ろした広さんが体ごと首を傾げる。

「それはちょっと…難しいと思う」

「やはり一朝一夕というわけにはいかない、と？」

「うん。というかそもそも星南たちも、どうやって見えてるのか説明できないと思う。暗黙知ってやつだ、ね」

今度は広さんが人差し指を立てる。『暗黙知』、聞いたことはあるが。

「…職人芸のようなもの…ということでしょうか」

「プロデューサー、知らないの？暗黙知」

意外、といわんばかりに身を乗り出す広さん。

「なんとなくは知っていますが、厳密な定義があるんですか？」

「昔、提唱者が書いた本\cite{polanyi1966ja}を読んだことがある。広せんせーが教えてあげようか？」

彼女はやけに私に勉強を教えたがる。一見すれば気の利いた申し出に、下心めいたものが感じられるのは気のせいだろうか。

「…わかりやすくお願いしますね？」

せめてもの命乞いに、腰に手を当てた広さんは自信に満ち溢れた表情で宣言する。

「まかせて」


\vskip\baselineskip
◆
\vskip\baselineskip
同室、黒板の前。

「"We can know more than we can tell"\cite[p4]{polanyi2009}だ、よ」


「……?」

「『我々は言葉にするより多くのことを知ることができる』\cite[18]{polanyi1966ja}、よ」

「日本語にすれば良いというものではないのですが」


悪い予感は的中した。思わず顔に当てた手のひらの、その指の隙間から得意げな広さんが見える。

「あれ、わからなかった？」

「やり直してみましょう」

「わかった。丁寧に説明してみる、ね」


こほん。と仕切り直す広さん。

「たとえば、プロデューサーは何万人といる人混みの中で、私の顔を見分けられる？」

「それは、もちろん」

「どうやって？」

「どう、と言いましても。広さんの顔は知っていますから。厳密に言えば目の色や、髪、眉の角度、とかでしょうか」

「普段からそんな細部を意識して私を見つけてるの？」

「意識は…していないですね」

「そう、意識していない。明示的な特徴を列挙できるからと言って、それを直接意識して顔を認知しているわけでもない。」

「つまり、何が言いたいんです？」

「プロデューサーが本当に自然体で顔を認知している時、それをどうやっているのかは、言語化できないはずだ、よ」

「ふうむ、そうなのかもしれません」

「あんまり納得してなさそう」

「普段から言語化を意識していないからといって、絶対に出来ないわけではないと思うのですが。それこそ、警察が事情聴取の内容から犯人そっくりの似顔絵を描く技術とかもありますし」

「そうだね。認知のメカニズムが解明され、それを適切に表現する方法が与えられたら、無意識下で行われていた認知を言葉で伝えることもできる。それでも、一つだけ変わらない事実がある」

「それは？」

「それらの方法が確立される前から、私たちはすでに人相を識別する方法を知っていた。つまりその時点では言葉にできるより多くのことを知っていた、ということだ、よ\cite[19]{polanyi1966ja}」

「なるほど」

「他によく例として挙げられるのは自転車の乗り方とか。自分がどう操縦しているか説明できる人はあまりいないだろうし、逆にいくら丁寧な説明を聞いても、それだけじゃ乗れるようにはならなさそうじゃない？」

「確かに。そう聞くとほとんどの身体的な知は言語化が難しそうですね」

「うん。だからそれらの知を身につけるには実際に体を動かしてみる他ない。練習あるのみだ、よ」

「ダンスも歌も、座学だけでマスターできたら広さんは今頃学園首席でしょうからね…」

「ふふ。そうじゃないから、たのしい」


上機嫌に微笑む広さん。

「ふむ…なるほど。つまり暗黙知とは、『まだ言語化されていない知識』ということでしょうか」

「そう聞こえちゃった？」

「今の説明ではそうとしか」

「今のはあくまで『言語化できていない知識がある』ということを示したかっただけであって、ポランニーによるとそれそのものが暗黙知ってわけじゃないらしい、よ」

「ほう」

「それは認知の構造にあるとポランニーは見たんだけど…」

「だけど？」

「…ふう」

「…座りましょうか」

◆

——同室、デスクにて。

淹れたコーヒーを半分ほど飲み切ったころに、広さんが口を開いた。

「さっきの話の続き、聞きたい？」

「ええ、ぜひ」

「ふふ…じゃあ、まずは実際に行われた実験の話。1956年にエリクセンとクーゼが行ったもの\footnote{ポランニー(高橋訳、2003)本文中(p24)では1958年とされているが、当該箇所の注釈では1956年の論文(Eriksen \& Kuethe, 1956)が示されている。}\cite{eriksen1956}なんだけど、被験者が会話の中で特定の『ショック語』に関連することを口にした時に、電気ショックを与える、といった内容」

「いきなり物騒ですね…」

「これがなんと、実験を繰り返すうちに被験者はショック語に関する言葉を口に出さなくなったんだって」

「パターンを見抜けるようになった、ということですか」

「無意識に、ね。被験者に尋ねたところ、本人はショックを回避していることには無自覚だったらしい、よ。この実験結果を見て、暗黙知の提唱者であるマイケル・ポランニーはその構造を提案した。まず暗黙知には二つの条件がある。この実験だとそれは第一条件であるショック語と、それに続く第二条件である電気ショック。被験者はこの二つを連結させて、ついにはショックを回避するまでに至ったけど、その連結は意識できるものには至らなかった。ポランニーはこれを、被験者が電気ショックに気を取られていたせいだと主張している\cite[26]{polanyi1966ja}、よ」

「まあ、電気ショックを認識せざるを得ないのは分かりますね」

「本来はここで二種類の認知が起きているはず。ショック語を察知することと、電気ショックを感じること。でも注意を払えているのは後者のみ。ポランニーはショック語（第一条件）について知るのは、電気ショック（第二条件）に注意を向け、その結果感知されたものを信じているからにすぎないと主張している\cite[27]{polanyi1966ja}、よ」


おもむろに手元のノートを開く広さん。

「第一条件から第二条件に意識が移動した結果、個別の諸要素が意味を持った包括的存在に統合されてゆく。これをポランニーは|包括（＝理解）《comprehend》と呼んだ\cite[33]{polanyi1966ja}、よ」


書かれた文字は、『ひろ』。

「ゲシュタルト崩壊\cite{gestaltzerfall}って、知ってる？」

「同じ文字を見つめ続けていると、その意味がわからなくなり、ただの記号に見えてくる現象ですよね」

「そう。これがまさに、個々の諸要素に注目しすぎた結果、それらの意味が消え、包括的な存在としての意味が失われている状態」

「確かに、普段文字を認識する際に細かい要素は意識しませんね」

「物事の諸要素を明瞭にすることは、それらの理解を深めてくれる時もあるけど、逆にそれらが持つ意味を拭い去ってしまう時もある。これらは認知だけじゃなくて、身体知にも現れることがある、よ」

「例えば？」

「一部のピアニストは、自分の指に注意を集中させると、一時的に演奏が出来なくなっちゃうらしい\cite[41]{polanyi1966ja}」

「なるほど。スポーツ等でも起こり得そうですね」

「うん。ポランニーはこうして諸要素を解き明かす明示性を絶対視することに疑問を投げかけていく」


「ふむ。とはいえ、明示性を追い求めることで得られるものもあると思いますがね」

「ポランニーも全否定はしていない、よ。実際にピアニストが自身の動きを明示的に研究することで、技能が改善される可能性も述べている。\cite[42]{polanyi1966ja}この人の主張はあくまで、諸要素の明示的統合と暗黙的統合はお互いに取って代わることは出来ないということ」

「まだ完全に納得は出来ませんね。その主張だと、『まだ明示的に記述出来ないだけ』というわけではなく、そもそも記述出来ない知識があることを示唆している気がします」

「察しがいいね、プロデューサー。ここからポランニーの主張はよりドラスティックになっていく、よ。まず、近代科学の部分批判が始まる」

「思い切りましたね」


「ポランニーによると、近代科学は主観的なものを排して、物事を形式化することを信条としている。\cite[44]{polanyi1966ja}たとえまだ形式化できないものがあっても、それは一時的なものに過ぎず、いずれは形式化されるであろうということ。ただ、ポランニーは暗黙的思考、言い換えれば主観的な認知が科学的方法の出発点として必要だと思った」

「その、暗黙的思考が必要という前提に根拠はあるのでしょうか」

「まだ形式化されていない知を形式化する為には、その存在を暗黙的に認知していないと、それを解き明かそうとは思えない、と言ったらどう？」

「ふうむ。それは確かにそうですが、形式化するにあたって、それを解き明かそうとしたという意図は必須ではないのでは？偶然発見される理論もあると思いますが」

「うん。このポランニーの主張は、普遍的なものじゃない。ただ、世の発見の中には明確に意図されたものもあって、その出発点は暗黙知だったとは言えるかも、ね」

「かなり主張が弱まりましたが、それは否定できないですね」

「加えておくと、仮に知の形式化が出来たとしても、その形式知の意味は暗黙的に理解され続けるとポランニーは主張している\cite[45]{polanyi1966ja}、よ。仮に、形式知としての理論があったとして、その理論を適用させる条件はまた別の理論になる。これを繰り返していくと無限後退になってしまう」

「その無限後退を断ち切るのが、包括（＝理解）であると？」

「ポランニーはそう考えた。暗黙的に認識することで形式知としての理論は不要になるから」

◆

「少し脱線したけど、ポランニーは科学的方法には暗黙的思考が必要不可欠だと主張しているから、それを元に研究を推し進める行為を孤独且つ個人的な行為だと思っている。他人に語れたら暗黙的ではないから。この発見についての予期とそれを探究しようとする行為をポランニーはコミットメントと呼び、その行為を形式化しようとすれば、行為そのものが台無しになると思った」\cite[52]{polanyi1966ja}

「先ほどの無限後退と同じように、形式化のキリがなさそうですね」

「これらの主張をもって、ポランニーは実証主義的な科学に一石を投じたということだ、ね」

「なるほど、彼のスタンスが分かってきました。ちなみに、暗黙的思考は人間にしか見られないんですかね？」

「最近では機械の暗黙的思考として、機械学習なんかが挙げられてるけど、まだ議論の余地がありそうだ、ね」

「実に興味深いですね」

◆

「ここまではポランニーが考える、暗黙知を用いた存在の知り方を説明した」

「知り方…そうですね」

「これで星南とおじいちゃんの能力も説明出来そうじゃない？」

「そういえば始まりはそのような話でしたね…。確かに、星南さん達がアイドルの諸要素を暗黙的に数値に統合する能力があるという解釈ができるかもしれません」

「うん。でも解釈はそこで終わらない、よ」

「ほほう」

「ポランニーは、暗黙知の構造が『知り方』だけじゃなくて『在り方』にまで適用できると思った」

「在り方…抽象的ですね」

「具体例を使って説明する前に、まずは暗黙知の構造についておさらいしてみよう」

「お願いします」

「存在を認識するには、まず具体的な諸要素を感知して、その感覚に依拠しながら、その存在に注意を向けていく。もし個々の諸要素に注意を移したなら、その機能は失われて、それまで注意を向けていた包括的存在を見失う」

「これが存在の在り方にも適用されると」

「うん。身近な例だと言語行動、つまり話すことの構造に適用できる、よ。そもそも私が今話しているのは、日本語の文法に則った文。つまり私の発言の諸要素は、文法だと言える。ただ、その発言の文法のみを検証し始めたら、その発言の意味は認識できなくなる」

「…英語の授業で、文法構造は理解できても、文の意味は理解出来なかったりしましたね。なるほど、諸要素に注目しすぎていたと」

「それか、ただの勉強不足？」

「…」

「こほん。ちなみにこの構造にはさらなる階層がある。文法の諸要素は単語、単語の諸要素は発音、発音の諸要素は喉や口の動き方とか、ね」

「この階層の下限や上限はあるのでしょうか」

「存在を証明することは難しそう。ただ、今の例だと文が文体論の諸要素になったり、文体論が文芸批評の諸要素になったりもする」

「それ以上は…ちょっと思い浮かびませんね」

「でも認識出来ないからって無いと決まったわけじゃない」

「下限に関しては、発話に関わる物理現象と考えることも出来そうですが」

「それもあくまで今の物理学で説明できる範囲での解釈でしかない。もしかしたらもっと細かい諸要素に分解できるかも」

「うーん、確かに下限も上限も見つかりそうにないですね」

「この階層構造は、諸要素で構成される下位層と、諸要素からなる包括的存在で構成される高位層の連なりと見ることが出来る」

「どの層も相対的には下位層でもあり、高位層でもあると」

「そう。そして各層は二つの規則によって制御される。その層自体に適用される規則と高位層を制御する規則」\cite[67]{polanyi1966ja}

「高位層の規則が下位層に及ぶ？」

「例えば発声は語彙に従って行われるし、語彙は文法によって並べられる。語彙を表さない発声は意味を成さないし、文法上で使い方が分からない語彙は存在し得ない」

「さらに、各層は下位層の規則を尊重しなければならない。例えば発声がでたらめになれば、単語は作れないし、存在しない単語を並べても意味のある文章にはならないから、ね」

「下位層が使えるパーツを決め、高位層がその使い方を決める、といった具合でしょうか」

「まさしくそう」

「そう言われると、使えるパーツを決める規則では、パーツ自体の使い方は決められないのも頷けますね」

「ポランニーはその高位から下位に向かって行われる制御を『境界制御の原理』\cite[73]{polanyi1966ja}と呼んだ、よ。そしてこの原理をもって理論的還元主義への批判が始まる」

「理論的還元主義とは？」

「全ての理論はより下位の理論に翻訳できるという考えだ、よ」

「なぜこれを批判しようと思ったのでしょうか」

「科学の多様な発展において必要だと思ったからじゃないかな。実際、還元主義を受け入れて全ての現象を物理学と化学に還元しようという動きは生物学など、様々な分野で見られている。でも『境界制御の原理』から分かる通り、それでは不十分。人体を物理学的、または化学的法則を用いて完璧に記述したとしても、それらの規則の視点からはそれが人体であることを判断することはできない。それを定義するのはさらに高位の法則、それこそ生物学上で定義されるもの」

「では何故、科学者達は還元主義的な研究を行うのでしょうか？」

「ここで区別したいのが、理論的還元主義と、方法論的還元主義。\cite{reductionism}ポランニーが批判しているのはあくまで前者。後者は方法論として、高位の事象を構成する諸要素を解き明かす行為。ポランニーはこれに関しては肯定的。諸要素を知ることでそれが高位に及ぼす制限を知ることが出来るから、ね」

「確かに、ピアニストの動作研究等に対して肯定的な意見もありましたね」

◆


「さて、『境界制御の原理』に基づき現象には意味の階層構造があると受け入れることで、新しく発生する概念がある、よ」

「それは？」

「『成功』と『失敗』」

「大きく出ましたね」


二本立てた指は二つの概念を表しているのだろうが、得意げな表情も相まってピースにしか見えない。

「ふふ。でも考えてみてほしい。アイドルの世界には成功と失敗の概念がある、でしょ？」

「何をもってアイドルが成功しているかを厳密に定義するのは難しいと思いますが…」

「ふむ。プロデューサー、厳密に定義出来ないからと言ってその概念が存在しないとは限らない」


意地悪な指摘だと自覚しているのであろう、いたずらっぽい表情を見せる広さん。

「暗黙的に認知している…ということでしょうか」

「プロデューサー。私、アイドルとして成功してる？」


話を合わせようと思った矢先、最も意地の悪い質問が飛んできた。

「…一番星がアイドルとして成功していないというのは無理がありますね」

「じゃあ一番星になってない千奈は？佑芽は？」

「彼女達も成功していると言えるでしょう」

「なんで？」

「…キリがなくなってきそうですね」

「ふふ。困ってるプロデューサー見るの、すき」


言葉通り、この上なく楽しそうに微笑んでいる広さん。

「認めましょう。少なくとも私は暗黙的にアイドルが成功しているか否かを認知できます」

「星南やおじいちゃんにアイドルの実力が見抜ける仕組みと変わらないんじゃない？」

「ふむ。そうかもしれませんが、丸め込まれた感が否めないですね」


肯定しきれなかったのは、意地悪された反感からだろうか。

「もっと成否が判別しやすいレベルまで落としてみよう、アイドルを構成する下位層においては定義が出来るものもある。例えば、振り付けを間違えたらダンスとしては失敗。ピッチを間違えたら歌としては失敗と言える」

「ふむ」

「でもダンスや歌を物理現象に還元してみた時、それが成功しているものなのか否かの区別は付けられない」

「存在論的還元主義の下では成功も失敗も存在し得ない、と」

「成功と失敗が存在しないと、どうなると思う？」


にまりと、広さんが笑う。
さっきの笑みとは違う。嫌味のない、わくわくしたような表情。

「？…どうなるんでしょうか」

「失敗のリスクが生まれない」


なるほど。

「そうなると――」


ピンと来た。


「「ままならなくなれない」」

「…死活問題ですね」

「でしょ？」

\begin{thebibliography}{99}

\bibitem{polanyi1966ja}
マイケル・ポランニー(高橋勇夫訳)(2003)『暗黙知の次元』筑摩書房(ちくま学芸文庫)

\bibitem{polanyi2009}
Polanyi, Michael (May 2009). The Tacit Dimension. Chicago: University of Chicago Press.

\bibitem{eriksen1956}
Eriksen, C. W., and Kuethe, J.L., "Avoidance Conditioning of Verbal Behavior Without Awareness:" A Paradigm of Repression," Journal of Abnormal and Social Psychology (Vol. 53, 1956). 203-09

\bibitem{gestaltzerfall}
「ゲシュタルト崩壊」
『Wikipedia』日本語版、
\url{https://ja.wikipedia.org/wiki/ゲシュタルト崩壊}
(閲覧日: 2026年9月22日)

\bibitem{reductionism}
「還元主義」
『Wikipedia』日本語版、
\url{https://ja.wikipedia.org/wiki/還元主義}
(閲覧日: 2026年9月22日)

\end{thebibliography}
\end{document}
